% SPDX-License-Identifier: Apache-2.0
% covers: Tracer
\datasheet{Tracer}{A ring of the last N things that happened, read out over AXI-Lite after the machine has stopped.}

\dsfacts{\code{lib/parts/src/tracer.rs}}{\code{txhdl\_parts::tracer::Tracer}}%
{\code{//docs:examples}}

\subsection*{Function}

\code{Tracer<N>} takes one 128-bit entry a cycle and keeps the last
\code{N} of them. What it holds is always the most recent window: when
the ring is full the oldest goes. A host reads the window out over
AXI-Lite, oldest first, by writing a cursor and reading the four words
of the entry the cursor names.

It is meant for what a machine was doing when it stopped, which is why
the readout is separate from the capture: a design raises \code{halt},
the buffer freezes, and the window is read at leisure afterwards over
the same bus the rest of the system uses.

An entry is 128 bits because a retirement is. A program counter, an
instruction, the register written and the word written come to a
hundred and one bits, and a width that is a whole number of bus words
is one the readout walks without shifting.

\subsection*{Parameters}

\begin{center}\footnotesize
\begin{tabular}{@{}lp{0.7\textwidth}@{}}
\toprule
Parameter & Meaning \\
\midrule
\code{N} & How many entries the ring holds. It is a power of two, since
the ring wraps by masking. The tables use 8, as the example and the
tests do. \\
\bottomrule
\end{tabular}
\end{center}

\subsection*{Register map}

Offsets from the peripheral's base. The buffer selects the word by bits
4 to 2 of the address and looks at no other bits, so the bridge in
front of it decides the base, and the seven words repeat every 32
bytes.

\dsregs{Tracer}

\code{CTRL\_RUN} and \code{CTRL\_FREEZE} are the masks of
\code{ctrl}'s two fields. \code{0x0c} is named by no register and reads
zero. A write to \code{count} or to a data word is taken and does
nothing.

\code{tracer::word(k)} gives the offset of word \code{k}, which is
\code{0x10 + 4k}.

\dstables{Tracer}

\subsection*{Behaviour}

\begin{itemize}
\item An entry is stored when the run bit is set, \code{take} is high,
and the buffer is not freezing. It goes into the slot \code{head}
names, \code{head} moves on by one and wraps by masking, and
\code{count} rises until it reaches \code{N} and then stays there.
\item What is held runs from the oldest to the slot before the head, so
a read at cursor \code{k} is the entry \code{k} after the oldest:
the slot is \code{(head - count + cursor)} masked to the ring.
\item The freeze is a level, not an edge. While \code{CTRL\_FREEZE} is
set and \code{halt} is high, the run bit is cleared, and it is cleared
again in every such cycle, so a host cannot start the buffer while
\code{halt} is still up. A host that reads \code{ctrl} back therefore
learns why the buffer stopped.
\item A write of the run bit while the buffer is stopped also zeroes
\code{count} and \code{head}: a host that starts the buffer starts a
new window rather than adding to the old one.
\item A read is answered on \code{bus\_r} in the cycle its address beat is
taken, with \code{OKAY}. It returns the registers as the last edge left
them, so a read in the same cycle as a write of the same word returns
the old value.
\item A write needs its address beat and its data beat in the same
cycle, and room on \code{bus\_b}; it is answered \code{OKAY} in that cycle.
The buffer holds no write address between cycles.
\item The cursor is a register, so a host writes it in one transaction
and reads the entry in the next.
\end{itemize}

\subsection*{Verification}

\begin{itemize}
\item Five tests in \code{lib/parts/src/tracer.rs}, in
\code{//lib/parts:parts\_test}, drive a \code{Tracer<8>} through a
\code{LiteHost}. \code{what\_is\_held\_is\_the\_last\_n\_entries} offers
twelve into a ring of eight and finds the fifth to the twelfth.
\code{fewer\_than\_a\_ring\_are}\allowbreak\code{\_held\_oldest\_first} offers three and
finds three, in order. \code{a\_halt\_freezes\_the\_window} offers three,
raises \code{halt}, offers six more, and finds the three and a cleared
run bit. \code{a\_start\_opens\_a\_new\_window} finds the count and the
window reset by a start. \code{a\_stopped\_buffer\_takes\_nothing} finds
nothing held when the buffer was never started.
\item \code{//lib/examples:ex\_tracer} puts it behind a
\code{LiteBridge<1, ..>} on an AXI4 link at \code{0x1000}, offers twelve
entries, raises \code{halt} and offers four more that the buffer must
refuse.
\item The build simulates \code{Tracer<8>}'s netlist against that run
under nvc and Verilator:
\code{//docs:tracer\_sim\_tracer\_tb\_test} and
\code{//docs:tracer\_vsim\_test}.
\end{itemize}

\subsection*{Used by}

Nothing in the tree yet. It is meant for the core's retirement wires,
which the board computes and throws away; issue 240 says so.

\subsection*{Limits}

\code{N} must be a power of two. The entry width is 128 bits and the
bus 32, and neither is a parameter. The buffer takes one entry a cycle,
since the ring has one write port. Byte strobes are ignored: a write of
\code{ctrl} or \code{cursor} takes the word's low bits whatever the
strobe says. Nothing is refused: an offset that names no word reads as
zero or as \code{ctrl}, and every read and every write is answered
\code{OKAY}. The cursor is not masked when it is written, only when the
slot is formed, so a cursor past \code{count} reads a slot that is in
the ring but not in the window. There is no reset port, no interrupt,
no timestamp and no flag that says the ring wrapped.
